\section{Discussion}
\label{sec:discussion}

\paragraph{Three routes to the same shadow price.} In ACE, the SCC is a ratio of costates, and the three families of methods studied in this series reach the costates in different ways. Deep equilibrium nets learn them as outputs that satisfy the residuals of the equilibrium conditions. Actor-critic methods obtain them as the gradient of a learned value function. SDDP computes them from the duals of convex stage problems, and Proposition~\ref{prop:gradient} shows that this is the same costate recursion that the equilibrium conditions encode, evaluated exactly at a visited state instead of enforced approximately over a distribution of states. Seen this way, SDDP is a form of approximate dynamic programming in which the value function is represented by cuts and improved with exact dual information; its accuracy in ACE comes from the fact that this representation contains the true value function.

\paragraph{Why ACE is easy for SDDP.} SDDP suits ACE unusually well: two decisions per period, convex stage problems, and an affine value function that a single cut represents exactly. Corollary~\ref{cor:two} turns this into convergence after one backward pass. This is a strength of the benchmark, since every result can be verified, but it also means that ACE does not test what makes SDDP hard in general: value functions that need many cuts, high-dimensional decisions, and many scenarios per stage. The two-iteration convergence and the absent warm-start gains should be read in that light. They say that neural SDDP cannot improve on an exact method that is already fast, not that warm starts are useless.

\paragraph{What the comparison with the deep equilibrium net shows.} Neural SDDP without iterations is about an order of magnitude more accurate than the pseudo-state deep equilibrium net on the same instances, at a quarter to a third of the training time. The comparison favours neural SDDP in two ways. First, its network is trained by supervised learning on exact solutions, which exist only because SDDP can solve every training instance; the deep equilibrium net learns from the equilibrium conditions alone and needs no solver. Second, the cut network has an easy target: a smooth map from three parameters and the period to six slopes and an intercept, the same at every state. The deep equilibrium net must represent policies and costates as functions of the full state, and its accuracy depends on where its simulated states fall. The fairer reading is that, when a convex solver can generate training data, learning the solution of the dual problem is much easier than learning the solution of the equilibrium conditions. Both approaches remain cheap to evaluate: the first-period SCC of a new instance needs only the predicted cut for period one and a single stage problem, so uncertainty quantification over many parameter draws is as feasible with neural SDDP as with the pseudo-state network.

\paragraph{Limits.} The results rest on three assumptions. The first is convexity. The stage problems are convex only because business-as-usual emissions follow an exogenous path; if they are tied to capital, as in DICE, emissions become a product of the abatement rate and a convex function of capital, the value function need not be concave in the state, and cuts are no longer valid upper bounds. Tipping points with thresholds or regime switches \parencite{LemoineTraeger2014,CaiLontzek2019} break convexity in the same way. SDDP can still be run as a heuristic in such models, but without its bounds, and neural SDDP inherits that. The second is stagewise independence: SDDP handles i.i.d.\ shocks directly, but persistent shocks or learning about uncertain parameters require adding states, which in turn makes the value functions curved. The third is the finite horizon: as Remark~\ref{rem:truncation} shows, carbon's shadow price converges only at the rate of discounting, so low discount rates require horizons of thousands of years, or a terminal value function. Finally, the cut network is trained with one seed per calibration, and I have not explored how many training instances are needed.

\paragraph{Where neural SDDP should pay off.} The case for neural SDDP is strongest exactly where ACE stops being a good benchmark: models whose value functions are curved enough to need many cuts, so that SDDP needs many iterations and good initial cuts save most of them. A natural next step is a convex extension of ACE that keeps a closed form for part of the problem, for example stochastic damages with a convex hazard, followed by models without closed forms, for which the methods validated here would themselves become the benchmark.
